\documentclass[11pt, openany, oneside]{book}
\usepackage{mathtools}
\usepackage{amsthm, amssymb}
\usepackage{graphicx}
\usepackage{subcaption}
\usepackage{microtype}
\usepackage[margin=1in]{geometry}
\usepackage[hidelinks]{hyperref}

\usepackage{fancyhdr}
\pagestyle{fancy}
\fancyhf{}
\fancyhead[R]{\thepage}
\fancyhead[L]{\rightmark}
\renewcommand{\headrulewidth}{0pt}

\newtheorem{thm}{Theorem}[section]
\newtheorem{prop}[thm]{Proposition}
\theoremstyle{definition}
\newtheorem{defn}[thm]{Definition}
\theoremstyle{remark}
\newtheorem{rmk}[thm]{Remark}

\newtheoremstyle{boldital}
  {\topsep}{\topsep}   % space above, below
  {\itshape}           % body font
  {}                   % indent
  {\bfseries}          % head font
  {.}                  % punctuation after head
  {.5em}               % space after head
  {}                   % head spec (default)

\theoremstyle{boldital}
\newtheorem*{replication}{Replication Pricing Scheme}
\newtheorem*{mart}{Martingale Pricing Scheme}

\numberwithin{equation}{section}

\title{Financial Foundations of SigAlg}
\author{John Myers}

\begin{document}

\frontmatter
\maketitle

\chapter*{Preface}
\addcontentsline{toc}{chapter}{Preface}
\markboth{Preface}{}
Text of the preface...

\tableofcontents

\mainmatter

\chapter{The binomial model and arbitrage pricing}

\section{The two-asset market and binomial model}

We suppose that the market contains two assets, one risky and the other risk-free, that we refer to as a \textit{stock} and \textit{bond}, respectively. We assume that the price process $S_t$ of the stock is stochastic, while the returns on the bond follow a per-period continuously compounded interest rate $r>0$. This means that $\$1$ invested in the bond will increase to a value of $e^{rt}$ over $t$ time periods. The rate $r$ is called the \textit{risk-free rate}. It will be convenient to write $R=e^r$, and call $R$ the \textit{gross risk-free return}.

We imagine an investor owns a two-asset portfolio consisting of positions in the stock and the bond. The portfolio is held over a discrete time interval beginning at $t=0$ and extending over $T$ time periods (so $t=0,1,\ldots,T$). At a given time $t$, we let $B_t$ be the value of the position in the bond, and we let $N_t$ be the number of units of stock held in the portfolio. The total value of the portfolio at time $t$ is then given by the process
	\[
	V_t \stackrel{\text{def}}{=} B_t + S_t N_t.
	\]
We allow $B_t$ and $N_t$ to take all real values, including negative ones, corresponding to short positions. 

As time ticks forward, we assume that the investor adjusts their positions in the stock and bond, so $B_t$ and $N_t$ are not constant processes. In fact, at a fixed time $t$, we assume that the investor will adjust \textit{both} their positions based on the price trajectory of the stock up to time $t$. This means that $B_t$ and $N_t$ are stochastic processes that are \textit{adapted} to the \textit{natural filtration} of the price process $S_t$: For each time $t$, the random variables $B_t$ and $N_t$ are measurable with respect to the $\sigma$-algebra
	\begin{equation}\label{eq:natural-filt}
	\mathcal{F}_t \stackrel{\text{def}}{=} \sigma(S_0,S_1,\ldots,S_t)
	\end{equation}
generated by the random variables $S_0,S_1,\ldots,S_t$.

So far we have placed no assumptions on the price process $S_t$ other than its stochasticity. Now, however, we introduce a simple model for the process. Let $u$, $d$, $p$ and $s_0$ be real numbers, with
	\[
	0 < d < 1 < u \quad \text{and} \quad 0 < p < 1,
	\]
and let $T$ be a positive integer. Let $(\Omega, \mathcal{F},P)$ be a probability space constructed as follows:
\begin{enumerate}
\item The sample space $\Omega$ consists of all $(T+1)$-tuples of the form
	\[
	\omega = (s_0, s_1,\ldots,s_T)
	\]
such that each quotient $s_{t+1} / s_t$ is either $u$ or $d$ (for $t=0,1,\ldots, T-1$). The number of times that $u$ appears as a quotient is called the number of \textit{up moves} in the sample point $\omega$, while the number of times that $d$ appears is called the number of \textit{down moves}.
\item The $\sigma$-algebra $\mathcal{F}$ is the power set on $\Omega$.
\item The probability measure $P$ is given by
	\begin{equation}\label{eq:real-world-measure}
	P \left( \{\omega\} \right) = p^n (1-p)^{T-n},
	\end{equation}
	where $n$ is the number of up moves in $\omega$.
\end{enumerate}
There are then $T+1$ coordinate projection random variables
	\begin{equation}\label{eq:binomial-proj}
	S_t: \Omega \to \mathbb{R}, \quad S_t\left( (s_0, s_1,\ldots,s_T)\right) = s_t,
	\end{equation}
with $S_0$ a constant random variable. 

\begin{defn}\label{defn:binomial}
The stochastic process $\{S_t\}_{t=0}^T$ is called a \textit{binomial model} with parameters $(u,d,p,s_0,T)$. The numbers $u$ and $d$ are called the \textit{up-factor} and \textit{down-factor}, $p$ and $P$ are called the \textit{real-world probability} and \textit{real-world measure}, the number $s_0$ is called the \textit{initial price}, and $T$ is called the \textit{time horizon} of the model.
\end{defn}

Note that the range of the random variable $S_t$ is the set $\{u^t d^{n-t}s_0 : n=0,1,\ldots,t\}$. By a simple counting argument, we have
	\begin{equation}\label{eq:binom-prob}
	P(X_t = u^t d^{n-t}s_0) = \binom{t}{n}p^t(1-p)^{t-n},
	\end{equation}
for each $n=0,1,\ldots,n$. The presence of the binomial coefficients is what gives the model its name.

Given a sample point $\omega\in \Omega$, its \textit{length-$t$ partial trajectory} is the sequence
	\[
	(S_0(\omega), S_1(\omega),\ldots,S_t(\omega)).
	\]
The next result identifies the atoms of the $\sigma$-algebras $\mathcal{F}_t$ in the natural filtration of a binomial model (see \eqref{eq:natural-filt}) and computes their probabilities in terms of these partial trajectories.\footnote{Recall that an \textit{atom} of a $\sigma$-algebra $\mathcal{F}$ is a minimal non-empty set in $\mathcal{F}$.} We leave the simple proofs to the reader.

\begin{prop}\label{prop:atom-prob}
Let $A$ be an atom of $\mathcal{F}_t$.
\begin{enumerate}
\item Two sample points $\omega$ and $\omega'$ are in $A$ if and only if their length-$t$ partial trajectories are equal.
\item We have $P(A) = p^n (1-p)^{t-n}$, where $n$ is the number of up moves in a length-$t$ partial trajectory of a sample point $\omega \in A$, which is independent of the choice of $\omega$ by \textnormal{(1)}.
\end{enumerate}
\end{prop}

Using this proposition, one easily proves the following measurability criterion.

\begin{prop}
Let $S_t$ be a binomial model with underlying sample space $\Omega$. A function $X: \Omega \to \mathbb{R}$ is $\mathcal{F}_t$-measurable if and only if it has the following property: Whenever $\omega$ and $\omega'$ are sample points in $\Omega$ such that
	\[
	(X_0(\omega), X_1(\omega), \ldots, X_t(\omega)) = (X_0(\omega'), X_1(\omega'), \ldots, X_t(\omega')),
	\] 
then $X(\omega) = X(\omega')$. \qedhere
\end{prop}

One of the most important and useful properties of the binomial model is given in the following result.

\begin{thm}
Every binomial model is a Markov process.
\begin{proof}
Let $S_t$ be a binomial model with time horizon $T$. It will suffice to show that
	\[
	P\left(S_{t+1} = s_{t+1} \mid S_1=s_1, \ldots,S_t=s_t\right) = P\left(S_{t+1} = s_{t+1} \mid S_t=s_t\right)
	\]
for all $t$ with $0 \leq t < T$, all $s_{t+1}\in \mathbb{R}$, and all vectors $(s_1,\ldots,s_t)$ of the form
	\[
	(s_1,\ldots,s_t) = (S_1(\omega),\ldots,S_t(\omega))
	\]
for some point $\omega$ in the sample space $\Omega$ of the model. To do this, note that
	\[
	P\left(S_{t+1} = s_{t+1} \mid S_1=s_1, \ldots,S_t=s_t\right) = \frac{P(S_1=s_1, \ldots,S_t=s_t,S_{t+1}=s_{t+1})}{P(S_1=s_1, \ldots,S_t=s_t)}
	\]
will be $0$ except in the cases that $s_{t+1}/s_t$ is either $u$ or $d$. In the first case, the probability in the numerator is $p$ times the probability in the denominator, while in the second case it is $1-p$ times the probability in the denominator. (Both of these computations use Proposition~\ref{prop:atom-prob}.) Thus the conditional probability is either $0$, $p$, or $1-p$.

On the other hand, we have
	\[
	P\left(S_{t+1} = s_{t+1} \mid S_t=s_t\right) = \frac{P(S_t=s_t, S_{t+1}=s_{t+1})}{P(S_t=s_t)}.
	\]
Again, this probability will be $0$ except in the cases that $s_{t+1}/s_t$ is either $u$ or $d$. In both of these latter cases, the probability of the denominator is
	\[
	P(S_t=s_t) = \binom{t}{n} p^n (1-p)^{t-n},
	\]
for some $n$ with $0\leq n \leq t$, from \eqref{eq:binom-prob}. But the same counting argument that establishes \eqref{eq:binom-prob} also shows that the probability of the numerator is either
	\[
	\binom{t}{n} p^{n+1} (1-p)^{t-n} \quad \text{or} \quad \binom{t}{n} p^{n} (1-p)^{t-n+1},
	\]
depending on whether $s_{t+1}/s_t$ is $u$ or $d$. The desired result then follows immediately.
\end{proof}
\end{thm}

\section[Arbitrage and martingales: The First Fundamental Theorem of Asset Pricing]%
        {Arbitrage and martingales: The First Fundamental Theorem of Asset Pricing%
         \sectionmark{The First Fundamental Theorem of Asset Pricing}}%
\sectionmark{The First Fundamental Theorem of Asset Pricing}

We continue with our two-asset market consisting of a risky stock and a risk-less bond, and throughout the rest of this chapter we assume that the price process $S_t$ of the stock is a binomial model with parameters $(u,d,p,s_0,T)$ as given in Definition~\ref{defn:binomial}. Let $R$ be the gross risk-free return of the bond.

We should expect that the potential gross return for the stock exceeds the guaranteed gross return for the bond since most investors are risk-averse and will demand a risk premium. This means that we should have $R < u$. However, this inequality also follows from requiring the market to be free of \textit{arbitrage}. In finance, this latter word generally refers to opportunities for a risk-free profit, and mathematical market models are designed in order to avoid them.

An arbitrage opportunity appears in the two-asset market if $R > u$. (We ignore the case that $R=u$.) For in this case the investor could short sell a unit of the stock at time $t=0$ and invest the cash into the bond; hence $N_0=-1$ and $B_0 = s_0$. Then, the initial value of the portfolio is
	\[
	V_0 = B_0 + s_0N_0 = s_0 - s_0 \cdot 1 =0,
	\]
while one time period later the value of portfolio is
	\[
	RB_0 + S_1N_0 = Rs_0 -S_1 \geq us_0 -us_0 > 0.
	\]
This portfolio has essentially created a sure profit out of nothing, which is an arbitrage.

So, in order to avoid arbitrage, we shall assume $R < u$. But $d<1$ and $R>1$, so we have the pair of inequalities
	\[
	d < R < u.
	\]
We shall refer to these inequalities as the \textit{no-arbitrage condition}. Note that the stock and bond were completely independent assets until we specified this condition. Indeed, this condition mixes $u$ and $d$, which are parameters of the stock, with $R$, which is a parameter of the bond. Therefore, because it expresses a link between the assets, the no-arbitrage condition is really a property of the market, and accordingly the market is said to be \textit{arbitrage-free} if the no-arbitrage condition holds.

The no-arbitrage condition guarantees that the number
	\[
	q \stackrel{\text{def}}{=} \frac{R-d}{u-d}
	\]
satisfies $0 < q < 1$. Using $q$, we may define a second probability measure $Q$ on the underlying measurable space $(\Omega, \mathcal{F})$ of the binomial model, by analogy with \eqref{eq:real-world-measure}, via
	\[
	Q(\{\omega\}) \stackrel{\text{def}}{=} q^n (1-q)^{T-n}
	\]
for each $\omega \in \Omega$, where $n$ is the number of up moves in $\omega$. Notice that this is the same formula \eqref{eq:real-world-measure} for the real-world measure $P$, but with $q$ in place of $p$.

This measure $Q$ is called the \textit{risk-neutral measure} of the model, terminology that will be explained in Section~\ref{sec:FTAP2}. Relative to $Q$, the probability of an up move is $q$, while the probability of a down move is
	\[
	1-q = \frac{u-R}{u-d}.
	\]
Notice that these two probabilities do not involve the real-world probability $p$ whatsoever, and so they have no relation to the real-world price dynamics of the stock. Rather, they are simply the (normalized) distances between the risk-free return $R$ and the up- and down-factors of the model: The risk-neutral probability of an up move is greater than that of a down move, or $q> 1-q$, when the gross risk-free return is closer to the up-factor than to the down-factor.

Though the probabilities they assign mean quite different things, the real-world and risk-neutral measures do have one feature in common, however: They agree on which events in $\Omega$ have probability $0$, which means the two measures are \textit{equivalent}.\footnote{In measure-theoretic language, two measures are \textit{equivalent} if each one is absolutely continuous with respect to the other.} Indeed, since $0<p<1$ and $0<q<1$, every singleton $\{\omega\}$ has strictly positive probability under both $P$ and $Q$, and hence so does every nonempty $A \subset \Omega$, being a union of such singletons. Thus $P(A)=0$ if and only if $A=\emptyset$, and similarly for $Q$, so the two measures are equivalent.

So, the absence of arbitrage in the two-asset market---with the stock price $S_t$ following a binomial model---implies the existence of the risk-neutral measure $Q$. As we will show in the fundamental result below, the risk-neutral measure has the special feature that the \textit{discounted price process} $\tilde{S}_t$, defined by 
	\[
	\tilde{S}_t = \frac{S_t}{R^t}
	\]
for each $t=0,1,\ldots,T$, is a martingale when expectations are computed relative to $Q$. For this reason, the risk-neutral measure is also called an \textit{equivalent martingale measure} (or \textit{EMM}).

\begin{thm}[First Fundamental Theorem of Asset Pricing --- Binomial Model]\label{thm:FTAP1-binomial}
Consider a two-asset market consisting of a risky stock and a risk-less bond, as described above, and suppose that the price process $S_t$ of the stock follows a binomial model with time horizon $T>0$ and natural filtration $\mathcal{F}_t$. Then the market is arbitrage-free if and only if there exists a measure $Q$, equivalent to the real-world measure $P$, for which the discounted price process of the stock is a martingale:
	\[
	\tilde{S}_t = E_Q ( \tilde{S}_r\mid \mathcal{F}_t),
	\]
for all $t$ and $r$ with $0 \leq t \leq r\leq T$.
\begin{proof}
Supposing that the market is arbitrage free, we already demonstrated that the risk-neutral measure $Q$ exists. We claim that the discounted price process is a martingale with respect to $Q$, and to prove this it will suffice to show
	\[
	RS_t(\omega) = E_Q\left(S_{t+1} \mid \mathcal{F}_t \right)(\omega)
	\]
for each $t=0,1,\ldots,T-1$ and for all $\omega\in \Omega$, where $\Omega$ is the sample space of the binomial model.

Fix $\omega$ and let $A$ be the atom of $\mathcal{F}_t$ containing it. Using the risk-neutral probability $q$, we may write
	\[
	Q(\{\omega\}) = \pi q^n (1-q)^{T-t-n}
	\]
for some integer $n$ and real number $\pi$ with $0\leq n \leq T-t$ and $0 < \pi \leq 1$. Let $A^u$ and $A^d$ consist of those $\eta \in A$ for which
	\[S_{t+1}(\eta) = uS_t(\omega) \quad \text{and} \quad S_{t+1}(\eta) = d S_t(\omega),
	\]
respectively. Then we have
	\[
	Q(A) = \pi, \quad Q(A^u) = \pi q, \quad \text{and} \quad Q(A^d) = \pi (1-q).
	\]
We now compute:
	\begin{align*}
	E\left(S_{t+1} \mid \mathcal{F}_t \right)(\omega) &= \frac{\sum_{\eta\in A} S_{t+1}(\eta)Q(\{\eta\})}{Q(A)} \\
	&= \frac{\sum_{\eta\in A^u} S_{t+1}(\eta)Q(\{\eta\}) + \sum_{\eta\in A^d} S_{t+1}(\eta)Q(\{\eta\})}{\pi} \\
	&= \left[ \frac{\sum_{\eta\in A^u} uQ(\{\eta\}) + \sum_{\eta\in A^d} dQ(\{\eta\})}{\pi} \right]S_t(\omega) \\
	&= \left[ \frac{\pi puq + \pi pd(1-q)}{\pi} \right]S_t(\omega) \\
	&= \left[ qu + (1-q)d\right]S_t(\omega) \\
	&= R S_t(\omega),
	\end{align*}
as desired.

For the converse, we suppose that $Q$ is an equivalent measure for which the discounted price process is a martingale. Let $\Omega^u$ and $\Omega^d$ be the two atoms of $\mathcal{F}_1 = \sigma(S_1)$, the first corresponding to those $\omega\in \Omega$ for which $S_1(\omega) = uS_0$, and the latter those $\omega$ for which $S_1(\omega) = dS_0$. Then, we compute:
	\begin{align*}
	RS_0 &= E_Q\left(S_1\right) \\
	&= \sum_{\omega\in \Omega} S_1(\omega) Q(\{\omega\}) \\
	&=\sum_{\omega\in \Omega^u} S_1(\omega) Q(\{\omega\}) + \sum_{\omega\in \Omega^d} S_1(\omega) Q(\{\omega\}) \\
	&= uS_0Q(\Omega^u) + dS_0Q(\Omega^d),
	\end{align*}
and so
	\begin{equation}\label{eq:unique-EMM}
	R = uQ(\Omega^u) + dQ(\Omega^d).
	\end{equation}
But $Q(\Omega^d) = 1 - Q(\Omega^u)$, from which it follows that $d < R < u$. 
\end{proof}
\end{thm}

\begin{rmk}\label{rmk:unique-EMM}
It follows from \eqref{eq:unique-EMM} that the EMM is unique. We do not include this uniqueness assertion in the First Fundamental Theorem, since it properly belongs in the Second Fundamental Theorem (Theorem~\ref{thm:FTAP2-binomial} below).
\end{rmk}


In fact, an EMM has another ``martingale property'' along with the one just identified in the First Fundamental Theorem. This one involves the value process $V_t$ of the two-asset portfolio held by our investor, defined by
	\begin{equation}\label{eq:value-process}
	V_t = B_t + S_tN_t,
	\end{equation}
for each $t=0,1,\ldots,T$. Just as the price process $S_t$ has its discounted version $\tilde{S}_t$, so too there is a discounted value process given by
	\[
	\tilde{V}_t = \frac{V_t}{R^t},
	\]
for each $t=0,1,\ldots,T$. The second martingale property possessed by $Q$ states that $\tilde{V}_t$ is a martingale when all expectations are computed relative to $Q$, provided that the two-asset portfolio is \textit{self-financing}. This latter terminology means that
	\begin{equation}\label{eq:self-finance}
	V_{t+1} = RB_t + S_{t+1}N_t.
	\end{equation}
for each $t=0,1,\ldots,T-1$.
	
We understand these equations as follows. Exactly at time $t$, the investor adjusts their bond and stock positions, resulting with $B_t$ invested in the bond and $N_t$ many units of stock. The value of the portfolio after the adjustment is then given by $V_t$ in \eqref{eq:value-process}. These positions are held over the interval $[t,t+1]$, resulting in a portfolio valued at
	\begin{equation}\label{eq:old-portfolio}
	RB_t + S_{t+1}N_t
	\end{equation}
at time $t+1$. The investor then adjusts their positions again, with $B_{t+1}$ invested in the bond and $N_{t+1}$ many units of stock, resulting in a portfolio valued at
	\[
	V_{t+1} = B_{t+1} + S_{t+1} N_{t+1}
	\]
at time $t+1$. The self-financing condition \eqref{eq:self-finance} requires that the total value $V_{t+1}$ of the new, adjusted portfolio at time $t+1$ is equal to the value \eqref{eq:old-portfolio} of the ``incoming'' portfolio held over the interval $[t,t+1]$. No capital infusions or withdrawals are required to build the new portfolio---in other words, it is self-financing.

Observe that the proof of the next martingale property of an EMM uses \textit{only} that the discounted price process $\tilde{S}_t$ is a martingale. It does not use anything specific to a binomial model.

\begin{thm}\label{thm:value-process-martingale}
Consider a two-asset market consisting of a risky stock and a risk-less bond, as described above, let $S_t$ be the price process of the stock with time horizon $T>0$, let $\mathcal{F}_t$ be the natural filtration of $S_t$, and let $V_t$ be the value process of a portfolio given by \eqref{eq:value-process}. Provided that there exists an equivalent martingale measure $Q$ and the value process $V_t$ is self-financing, the discounted value process is a martingale:
	\[
	\tilde{V}_t = E_Q ( \tilde{V}_r \mid \mathcal{F}_t )
	\]
for $0 \leq t \leq r \leq T$.

\begin{proof}
It will suffice to prove
	\[
	RV_t = E_Q\left( V_{t+1} \mid \mathcal{F}_t \right)
	\]
for each $t=0,1,\ldots,T-1$. Since the portfolio is self-financing, we have
	\[
	V_{t+1} = R B_t + S_{t+1} N_t.
	\]
We then compute:
	\begin{align*}
	E_Q\left( V_{t+1} \mid \mathcal{F}_t \right) &= RE_Q(B_t \mid \mathcal{F}_t) + E_Q(S_{t+1}N_t \mid \mathcal{F}_t) \\
	&= RB_t + E_Q(S_{t+1} \mid \mathcal{F}_t) N_t \\
	&= RB_t + RS_t N_t \\
	&= R V_t,
	\end{align*}
where we used $\mathcal{F}_t$-measurability of $B_t$ and $N_t$ in going from the first line to the second, and that the discounted process $\tilde{S}_t$ is a martingale in going from the second to the third.
\end{proof}
\end{thm}


\section[Replicating portfolios: The Second Fundamental Theorem of Asset Pricing]%
        {Replicating portfolios: The Second Fundamental Theorem of Asset Pricing%
         \sectionmark{The Second Fundamental Theorem of Asset Pricing}}%
\sectionmark{The Second Fundamental Theorem of Asset Pricing}\label{sec:FTAP2}

By way of motivation, suppose that $S_t$ is the price process of a stock, with initial price $S_0$ at time $t=0$. At this initial time, two parties agree on a contract, called a \textit{European call option}, that gives one of them (the \textit{option holder}, or \textit{buyer}) the right, but not the obligation, to purchase a number of units of the stock from the other party (the \textit{option writer}, or \textit{seller}) at some future time $t=T$ at some price $\$K$ per stock. All of these parameters---the \textit{maturity} $T$ and the \textit{strike price} $\$K$, as well as the number of stocks to be purchased---are set in the contract at the initial time $t=0$.

The option holder will choose to exercise their right to buy the stocks provided that the actual market price $\$S_T$ of the stock at maturity $t=T$ (the \textit{spot price} at $t=T$) is greater than the strike price $\$K$. For in this case, the buyer can immediately turn around and sell the stocks on the market at $\$S_T$ per share, pocketing $\$(S_T-K)$ per share in profit. If instead the strike exceeds the spot price, the option holder will choose \textit{not} to exercise their right to buy, and the contract will expire at $t=T$ without any exchange.

An option contract costs money. At the initial time $t=0$, the holder pays a premium to the writer to purchase the rights afforded to them in the contract. \textit{The basic problem is to identify a fair value for this premium.} Basic finance tells us that the price of a financial asset should be the sum of all present values of future cash flows, but if the asset is an option contract like the one described above, we are immediately presented with two problems:
\begin{enumerate}
\item \textit{The future payoff is random}, depending (as it does) on the future spot price $S_T$ of the stock.
\item \textit{There is no obvious choice for the discount factor} with which to compute present values of future cash flows.
\end{enumerate}

One approach that solves these problems is to embed the stock into a market along with a risk-less asset---a market identical to the two-asset one we have been discussing so far---and to appeal to an arbitrage argument. Indeed, imagine that it is possible to construct a portfolio, including positions in a risk-less bond and the stock, in such a way that the final value\footnote{The notation is the same as the previous sections: $R$ is the risk-free return of the bond, $B_t$ is the value of the position in the bond at time $t$, $N_t$ is the number of stocks held at time $t$, and $V_t$ is the total value of the portfolio at time $t$.}
	\[
	V_T = RB_{T-1} + S_{T} N_{T-1}
	\]
of the portfolio held over the time interval $[T-1, T]$ coincides exactly with the payoff of the call option at maturity $t=T$. Then, provided that the portfolio is self-financing, an investor who purchases such a portfolio makes a \textit{single} outlay of capital at time $t=0$ to purchase it, and at maturity $t=T$ they have a portfolio whose value perfectly replicates the payoff of the call option. If the purchase price of the call option were not equal to the purchase price of the replicating portfolio, then an arbitrage occurs: The investor would only have to short sell the more expensive asset at time $t=0$ (either the call option or the replicating portfolio), buy the less expensive one with the cash from the sale, pocketing the difference as profit, and then cover the short position at time $t=T$ with the payoff.

This leads us to the following pricing scheme.

\begin{replication}[European call options]
Construct, if possible, a self-financing portfolio that perfectly replicates the payoff of the call option. Then, provided that such a portfolio is unique, the purchase price of the call option is taken to be the purchase price of the replicating portfolio.
\end{replication}

This scheme solves the two problems identified above, as long as an equivalent martingale measure $Q$ exists. Indeed, if $V_t$ is the value process of the unique self-financing portfolio that replicates the payoff of the call option, Theorem~\ref{thm:value-process-martingale} tells us that
\[
	V_0 = \frac{1}{R^T}E_Q(V_T),
\]
where the initial price $V_0$ of the replicating portfolio is (by our arbitrage pricing scheme) the purchase price of the call option and $V_T$ is (by definition of \textit{replication}) its payoff at maturity $t=T$. So, the price of the call option is indeed the present value of its expected future cash flow, where the expectation is computed relative to an EMM and the discount factor is the risk-free return.

This leads us to a second pricing scheme.

\begin{mart}[European call options]
Let $X_t$ be the price process of the call option with maturity $t=T$, so that $X_0$ is its initial price and $X_T$ is its payoff at maturity, and let $R$ be the risk-free return of the bond. Construct, if possible, an equivalent martingale measure $Q$. Then, provided that $Q$ is unique, the purchase price of the call option is given by
\[
	X_0 = \frac{1}{R^T}E_Q(X_T).
\]
\end{mart}

As we indicated, Theorem~\ref{thm:value-process-martingale} provides the bridge between the two pricing schemes, showing that they give the same price for a call option, contingent on the existence and uniqueness of replicating portfolios (for the Arbitrage Pricing Scheme) and equivalent martingale measures (for the Martingale Pricing Scheme). Whether these objects exist and are unique depends, in turn, on the underlying model for the price process $S_t$ of the stock. If this process is a binomial model, then the First Fundamental Theorem of Asset Pricing (Theorem~\ref{thm:FTAP1-binomial}) shows that an EMM exists, provided that the market is arbitrage free. Note also that the no-arbitrage condition provides uniqueness of the arbitrage price of a call option given as the initial value of a replicating portfolio—indeed, if there were two replicating portfolios with different initial prices, then an arbitrage opportunity appears.

So, as long as $S_t$ is a binomial model, the First Fundamental Theorem answers the uniqueness question of the Arbitrage Pricing Scheme and the existence question of the Martingale Pricing Scheme. The main result of this section—the Second Fundamental Theorem of Asset Pricing—provides answers to the remaining questions, the existence question in the Arbitrage Pricing Scheme and the uniqueness question in the Martingale Pricing Scheme. (See Table~\ref{tab:ftap}.)

\begin{table}[t]
\centering
\begin{tabular}{lcc}
\hline
 & \textbf{Existence} & \textbf{Uniqueness} \\
\hline
\textbf{Arbitrage Pricing Scheme}  & FTAP2 & FTAP1 \\
\textbf{Martingale Pricing Scheme} & FTAP1  & FTAP2 \\
\hline
\end{tabular}
\caption{Summary of the Fundamental Theorems of Asset Pricing (FTAP)}
\label{tab:ftap}
\end{table}

We will prove the Second Fundamental Theorem for binomial models. In fact, we will prove it for contracts more general than the European call options we have been discussing so far. These contracts are defined as follows.

\begin{defn}\label{defn:contingent-claim}
Let $S_t$ be the price process of the stock, indexed $t=0,1,\ldots,T$. Then a \textit{contingent claim}, with \textit{maturity} $T$, is an $S_t$-adapted process $\Phi_t$.

If $\Phi_t=0$ for all $t$ with $0 \leq t < T$, then the claim is called a \textit{non-early-exercise} claim, and the claim is said to \textit{not allow early exercise}. Otherwise, it is called an \textit{early-exercise} claim, and the claim is said to \textit{allow early exercise}.

The stock is called the \textit{underlying asset} of the claim, or simply the \textit{underlying}.
\end{defn}

The primary types of contingent claims discussed in these notes are defined as follows.

\begin{defn}\label{defn:options}
Let $S_t$ be the price process of the stock, indexed $t=0,1,\ldots,T$, and let $K$ be a positive real number, called a \textit{strike price}.
\begin{enumerate}
\item The \textit{European call option}, with strike $K$ and maturity $T$, is given by
	\[
	\Phi^\text{Euro-call}_t = \begin{cases} 0 & : 0 \leq t < T, \\ S_T-K &: t=T.\end{cases}
	\]
The \textit{European put option}, with strike $K$ and maturity $T$, is given by
	\[
	\Phi^\text{Euro-put}_t = \begin{cases} 0 & : 0 \leq t < T, \\ K-S_T &: t=T.\end{cases}
	\]
\item Write $\overline{S}$ for the average stock price, given by
	\[
	\overline{S} = \frac{S_0+S_1+\cdots + S_T}{T+1}.
	\]
Then the \textit{Asian call option}, with strike $K$ and maturity $T$, is given by
	\[
	\Phi^\text{Asian-call}_t = \begin{cases} 0 & : 0 \leq t < T, \\ \overline{S}-K &: t=T.\end{cases}
	\]
The \textit{Asian put option}, with strike $K$ and maturity $T$, is given by
	\[
	\Phi^\text{Asian-put}_t = \begin{cases} 0 & : 0 \leq t < T, \\  K-\overline{S} &: t=T.\end{cases}
	\]
\item The \textit{American call} and \textit{put options}, both with strike $K$ and maturity $T$, are given by
	\[
	\Phi^\text{American-call}_t = S_t-K \quad \text{and} \quad \Phi^\text{American-put}_t = K-S_t
	\]
for $0\leq t \leq T$.
\end{enumerate}
\end{defn}

While the process $\Phi_t$ itself is called a contingent claim, we shall also often call it the \textit{payoff process} of the contingent claim, calling back to the ``real-world'' scenario described at the beginning of this section. In this scenario, we imagine that a buyer purchases the contingent claim from a counterparty, which gives them the rights to \textit{one} of the future cash flows $\Phi_0,\Phi_1,\ldots,\Phi_T$. The time $t$ of the cash flow $\Phi_t$ selected by the buyer is called the \textit{exercise time}. Crucially, while the maturity $T$ is set and agreed upon in the purchase contract for the contingent claim at time $t=0$, the exercise time is not: At each time $t$ up to maturity, the buyer may choose to exercise their rights immediately and collect the cash flow $\Phi_t$, or defer exercise to a future time. If exercised, the contingent claim then expires; if the buyer does not exercise at maturity, the claim also expires. In the latter case, it will be convenient to say that the holder exercises at time $t=\infty$, so that the set of possible exercise times is $\{0,1,\ldots,T,\infty\}$.

Obviously, the holder wants to maximize their profit. They should then avoid exercising at times $t$ when $\Phi_t=0$ (since they collect nothing) or $\Phi_t<0$ (which we interpret as the holder owing the seller money). This explains the \textit{early exercise} terminology in Definition~\ref{defn:contingent-claim}: A holder of a non-early-exercise contingent claim would never exercise at any time besides $t=T$ (maturity) or $t=\infty$ (never).

\begin{thm}[Second Fundamental Theorem of Asset Pricing --- Binomial Model]\label{thm:FTAP2-binomial}
Consider a two-asset market consisting of a risky stock and a risk-less bond, as described above, and suppose that the price process $S_t$ of the stock follows a binomial model. Provided that the market is arbitrage free, we have that:
\begin{enumerate}
\item Every non-early-exercise contingent claim has a unique self-financing replicating portfolio.
\item The equivalent martingale measure is unique.
\end{enumerate}
\begin{proof}
The second statement was already established in the proof of the First Fundamental Theorem; see Remark~\ref{rmk:unique-EMM}.

The method of proof for the first statement is a familiar ``backwards induction,'' as it is known in the literature. Given a contingent claim $\Phi$ with maturity $T>0$, the idea is to construct the desired replicating portfolio processes $B_t$ and $N_t$ backwards
	\[
	B_{T-1},B_{T-2} \ldots, B_{t+1}, \quad \text{and} \quad N_{T-1},N_{T-2} \ldots, N_{t+1},
	\]
beginning with $B_{T-1}$ and $N_{T-1}$, and going back to some time $t+1$ with $0 \leq t \leq T-2$. The base case is a construction of $B_{T-1}$ and $N_{T-1}$, while the inductive step is to construct $B_t$ and $N_t$ from $B_{t+1}$ and $N_{t+1}$. It turns out that the construction for the base case is the same as the construction for the inductive step, so we go straight to the proof of the latter.

Assuming $B_{t+1}$ and $N_{t+1}$ exist, we have the portfolio value
	\[
	V_{t+1} = B_{t+1} + S_{t+1}N_{t+1}.
	\]
But the portfolio is supposed to be self-financing, and so we must also have
	\begin{equation}\label{eq:self-financing}
	V_{t+1} = R B_t + S_{t+1}N_t.
	\end{equation}
Note that this is an equation in the two unknowns $B_t$ and $N_t$. We shall prove that it has a unique solution.

Since both $B_t$ and $N_t$ are to be $\mathcal{F}_t$-measurable (where $\mathcal{F}_t$ is the $\sigma$-algebra generated by $S_0,S_1,\ldots,S_t$) we need only define these random variables on the atoms of this $\sigma$-algebra. Let $A$ be such an atom, and let $\omega \in A$. The atom contains (at least) two trajectories, $\omega_u$ and $\omega_d$, with the properties that
	\[
S_{t+1}(\omega_u) = u S_{t}(\omega) \quad \text{and} \quad S_{t+1}(\omega_d) = d S_{t}(\omega),
	\]
where $u$ and $d$ are the up- and down-factors of the model. Then, we have
	\[
	V_{t+1}(\omega_u) = R B_{t}(\omega) + uS_{t}(\omega) N_{t}(\omega)
	\]
and
	\[
	V_{t+1}(\omega_d) = R B_{t}(\omega) + dS_{t}(\omega) N_{t}(\omega)
	\]
from \eqref{eq:self-financing} and the equations
	\[B_t(\omega) = B_t(\omega_u) = B_t(\omega_d) \quad \text{and} \quad N_t(\omega) = N_t(\omega_u) = N_t(\omega_d).
	\]
We then have the linear system
	\begin{equation}\label{eq:linear-sys}
	\begin{bmatrix}
	R & uS_{t}(\omega) \\
	R & dS_{t}(\omega)
	\end{bmatrix}
	\begin{bmatrix} B_{t}(\omega) \\ N_{t}(\omega) \end{bmatrix} =
	\begin{bmatrix}
	V_{t+1}(\omega_u) \\
	V_{t+1}(\omega_d)
	\end{bmatrix}
	\end{equation}
in the unknowns $B_{t}(\omega)$ and $N_{t}(\omega)$. The determinant of the coefficient matrix is $RS_t(\omega)(d-u)$, which is never $0$, and hence a unique solution exists. This completes the proof of the inductive step.

Now, we return to the base case, demonstrating how to construct $B_{T-1}$ and $N_{T-1}$. The portfolio is supposed to replicate the contingent claim, and hence we must have
	\[
	\Phi^+_T = RB_{T-1} + S_T N_{T-1}.
	\]
This is an equation in the two unknowns $B_{T-1}$ and $N_{T-1}$ which is of the same form as \eqref{eq:self-financing}. The same argument as above shows that it has a unique solution. This completes the proof of the first claim in the theorem.
\end{proof}
\end{thm}


\newpage
Indeed, we imagine that these latter contracts give the holder the right, but not the obligation, to \textit{exercise} at any time $t$ and collect $\$\Phi_t$ from the writer. We interpret $\Phi_t<0$ to mean that the \textit{holder owes the writer} an amount $\$(-\Phi_t)$, obviously an undesirable outcome from the holder's perspective, and so the holder would never choose to exercise at these times. Once the holder decides to exercise, the contract expires. The holder also has the option \textit{never} to exercise and to let the contract mature worthless---in this last case, we say that the holder ``exercises'' at time $t=\infty$.

Once the holder has a contingent claim, they need to solve an optimization problem: With the goal to maximize their profit, when  should they exercise? For the holder of a non-early-exercise claim, the optimal exercise times are obvious since the only possible time to turn a positive profit comes at maturity. Therefore, the holder should exercise at $t=T$, if $\Phi_T>0$. Otherwise, if $\Phi_T \leq 0$, the owner should let the contract mature worthless. Optimal exercise times for early-exercise claims take much more work to identify, which we save for Section~\ref{sec:american}.

These contingent claims cost money. At the initial time $t=0$, the holder pays a premium to the writer to purchase the rights afforded to them in the contract. The basic problem is to identify a fair value for this premium.

Basic finance tells us that the price of an asset should be the sum of all present values of future cash flows, but if the asset is a contingent claim, we are immediately presented with three problems:
\begin{enumerate}
\item The holder can exercise at any time and collect their payoff. But this time will depend on the price trajectory of the stock, and hence \textit{the exercise time is random.}
\item Even if we knew the exercise time, \textit{the payoff itself is random}.
\item Even if we knew both the exercise time and the exact payoff, \textit{there is no obvious choice for the discount factor} with which to compute its present value.
\end{enumerate}
Of course, as we mentioned, the first is really only a problem for early-exercise claims. For European and Asian options and others like them, which do not allow early exercise, we can remove $\infty$ from the set of allowed exercise times and force exercise at maturity $t=T$, while still keeping the same cash flows, by replacing $\Phi_T$ with
	\begin{equation}\label{eq:zero-trick}
	\Phi_T^+ \stackrel{\text{def}}{=} \max{\left(\Phi_T, 0\right)}.
	\end{equation}
For a European call, this means that we consider
	\[
	\Phi_T^+ = \max{\left(S_T - K, 0 \right)}.
	\]
If $S_T< K$, then exercise at maturity $t=T$ results in the holder collecting $\$0$ from their counterparty, which is the same as the original European call with exercise time $t=\infty$.

\textit{In this first section on contingent claims, we shall only consider pricing non-early-exercise claims}, and leave the more complicated ones to a later section.

Now, one approach that attempts to overcome the second and third problems identified above is to embed the underlying asset into a market along with a risk-less asset---a market identical to the two-asset one we have been discussing so far---and to appeal to a no-arbitrage argument. Indeed, consider a (non-early-exercise) contingent claim $\Phi_T$ on our risky stock. Whatever its fair price $X_0$ at time $t=0$ turns out to be, imagine that it is possible to construct a portfolio, including positions in the bond and the stock, in such a way that the final value
	\[
	RB_{T-1} + S_{T} \Delta_{T-1}
	\]
of the portfolio held from $[T-1, T]$ coincides exactly with the final payoff $\Phi_T^+$. Then, provided that the portfolio is also self-financing, an investor who purchases such a portfolio makes a \textit{single} outlay of capital at time $t=0$ to purchase the portfolio, and ends at time $t=T$ with a portfolio whose value perfectly replicates the claim. If the purchase price $X_0$ of the claim was not equal to the purchase price $V_0$ of the replicating portfolio, then an arbitrage occurs: The investor would only have to short sell the more expensive asset at time $t=0$ (either the claim or the replicating portfolio), buy the less expensive one with the cash from the sale, pocketing the difference as profit, and then cover the short position at time $t=T$ with the payoff.

We are thus led to the following no-arbitrage pricing scheme: Given a (non-early-exercise) contingent claim $\Phi_T$, construct a self-financing portfolio that perfectly replicates the payoff $\Phi_T^+$ of the claim. Then, provided that such a portfolio is unique, the purchase price $X_0$ of the claim is taken to be the initial value $V_0$ of the replicating portfolio.

In order for this pricing strategy to work, we require both \textit{existence} and \textit{uniqueness} of these replicating portfolios. In the case that the price process for the underlying stock is a binomial model, these requirements are easily handled:

\begin{thm}[Second Fundamental Theorem of Asset Pricing --- Binomial Model]\label{thm:FTAP2-binomial}
Consider a two-asset market consisting of a risky stock and a risk-less bond, as described above, and suppose that the price process $S_t$ of the stock follows a binomial model. Provided that the market is arbitrage free, we have:
\begin{enumerate}
\item Every non-early-exercise contingent claim has a unique self-financing replicating portfolio.
\item The equivalent martingale measure is unique.
\end{enumerate}
\begin{proof}
The method of proof for the first statement is a familiar ``backwards induction,'' as it is known in the literature. Given a contingent claim $\Phi$ with maturity $T>0$, the idea is to construct the desired replicating portfolio processes $B_t$ and $N_t$ backwards
	\[
	B_{T-1},B_{T-2} \ldots, B_{t+1}, \quad \text{and} \quad N_{T-1},N_{T-2} \ldots, N_{t+1},
	\]
beginning with $B_{T-1}$ and $N_{T-1}$, and going back to some time $t+1$ with $0 \leq t \leq T-2$. The base case is a construction of $B_{T-1}$ and $N_{T-1}$, while the inductive step is to construct $B_t$ and $N_t$ from $B_{t+1}$ and $N_{t+1}$. It turns out that the construction for the base case is the same as the construction for the inductive step, so we go straight to the proof of the latter.

Assuming $B_{t+1}$ and $N_{t+1}$ exist, we have the portfolio value
	\[
	V_{t+1} = B_{t+1} + S_{t+1}N_{t+1}.
	\]
But the portfolio is supposed to be self-financing, and so we must also have
	\begin{equation}\label{eq:self-financing}
	V_{t+1} = R B_t + S_{t+1}N_t.
	\end{equation}
Note that this is an equation in the two unknowns $B_t$ and $N_t$. We shall prove that it has a unique solution.

Since both $B_t$ and $N_t$ are to be $\mathcal{F}_t$-measurable (where $\mathcal{F}_t$ is the $\sigma$-algebra generated by $S_0,S_1,\ldots,S_t$) we need only define these random variables on the atoms of this $\sigma$-algebra. Let $A$ be such an atom, and let $\omega \in A$. The atom contains (at least) two trajectories, $\omega_u$ and $\omega_d$, with the properties that
	\[
S_{t+1}(\omega_u) = u S_{t}(\omega) \quad \text{and} \quad S_{t+1}(\omega_d) = d S_{t}(\omega),
	\]
where $u$ and $d$ are the up- and down-factors of the model. Then, we have
	\[
	V_{t+1}(\omega_u) = R B_{t}(\omega) + uS_{t}(\omega) N_{t}(\omega)
	\]
and
	\[
	V_{t+1}(\omega_d) = R B_{t}(\omega) + dS_{t}(\omega) N_{t}(\omega)
	\]
from \eqref{eq:self-financing} and the equations
	\[B_t(\omega) = B_t(\omega_u) = B_t(\omega_d) \quad \text{and} \quad N_t(\omega) = N_t(\omega_u) = N_t(\omega_d).
	\]
We then have the linear system
	\begin{equation}\label{eq:linear-sys}
	\begin{bmatrix}
	R & uS_{t}(\omega) \\
	R & dS_{t}(\omega)
	\end{bmatrix}
	\begin{bmatrix} B_{t}(\omega) \\ N_{t}(\omega) \end{bmatrix} =
	\begin{bmatrix}
	V_{t+1}(\omega_u) \\
	V_{t+1}(\omega_d)
	\end{bmatrix}
	\end{equation}
in the unknowns $B_{t}(\omega)$ and $N_{t}(\omega)$. The determinant of the coefficient matrix is $RS_t(\omega)(d-u)$, which is never $0$, and hence a unique solution exists. This completes the proof of the inductive step.

Now, we return to the base case, demonstrating how to construct $B_{T-1}$ and $N_{T-1}$. The portfolio is supposed to replicate the contingent claim, and hence we must have
	\[
	\Phi^+_T = RB_{T-1} + S_T N_{T-1}.
	\]
This is an equation in the two unknowns $B_{T-1}$ and $N_{T-1}$ which is of the same form as \eqref{eq:self-financing}. The same argument as above shows that it has a unique solution. This completes the proof of the first claim in the theorem.

Finally, let $Q$ be an equivalent martingale measure. Since the discounted price process of the stock is a martingale, we must have $RS_0 = E_Q(S_1)$. If $A^u$ is the atom of $\mathcal{F}_1 = \sigma(S_1)$ consisting of those $\omega\in \Omega$ whose length-$1$ partial trajectories are equal to $(S_0,uS_0)$, and $A^d$ is the same atom but with $d$ in place of $u$, then we have
	\[
	E_Q(S_1) = \sum_{\omega\in A^u} S_1(\omega)Q(\{\omega\}) + \sum_{\omega\in A^d} S_1(\omega)Q(\{\omega\}) = uS_0 Q(A^u) + dS_0Q(A^d) = uS_0q + dS_0(1-q).
	\]
Thus $R = uq + d(1-q)$, which shows that $q$ is unique.
\end{proof}
\end{thm}

Provided that self-financing, replicating portfolios exist and are unique, the no-arbitrage argument was used to show that the initial price $X_0$ of a non-early-exercise contingent claim $\Phi_T$ on the stock coincides with the initial value $V_0$ of its replicating portfolio. It is natural to extend this to all times $t=0,1,\ldots,T$, defining the entire price process $X_t$ of the claim, by setting $X_t=V_t$, where we take $V_T = \Phi_T^+$. In the case that the price process of the stock is a binomial model, the no-arbitrage condition also provides the existence of a unique equivalent martingale measure $Q$, by the First Fundamental Theorem. From Theorem~\ref{thm:value-process-martingale}, we then get concise formulas for the no-arbitrage price process of the claim:
	\begin{equation}\label{eq:martingale-pricing}
	X_t = \frac{1}{R^{r-t}} E_Q\left( X_r \mid \mathcal{F}_t\right),
	\end{equation}
for $0 \leq t \leq r \leq T$. In particular, the initial price is given by
	\begin{equation}\label{eq:expectation-pricing}
	X_0 = \frac{1}{R^T} E_Q\left( \Phi_T^+ \right).
	\end{equation}
This latter formula is especially enlightening in view of the discussion at the beginning of this section, and how it navigates the second two problems listed there:
\begin{enumerate}
\item[(2)] The problem of not knowing the future payoff of the contingent claim with certainty, because it is a function of a random variable, is circumvented by considering the \textit{expected} future payoff. However, one must take care to note that \textit{the expected value is computed relative to the risk-neutral measure $Q$}, and not the real-world measure $P$.
\item[(3)] The problem of not knowing the correct discount factor is solved by using the risk-free rate. This is of interest since, in the real world, the correct discount factor would take into account the risk profile of the claim and the stock; however, in this formula, the risk-neutral measure $Q$ absorbs and reflects this risk, enabling one to discount at the risk-free rate.
\end{enumerate}
Since the formula \eqref{eq:expectation-pricing} prices the claim in terms of its expected payoff only, it is a \textit{risk-neutral} formula, which inspires the name of the risk-neutral measure $Q$.\footnote{An analogy is helpful. Suppose a player is selecting one of two games to play. In the first, the player wins or loses $\$100$ with equal probability $0.5$, while in the second, they win or lose $\$10$, also with equal probability. A \textit{risk-seeking} player would choose the first game and a shot at higher winnings, at the expense of risking higher losses, while a \textit{risk-averse} player would choose the latter, less risky game. But a \textit{risk-neutral} player would only consider the expected payoff of each game, which is $\$0$ for both, and would thus not express a preference for one game over the other. (Admittedly, the analogy with the risk-neutral measure is not perfect, because $Q$ \textit{does} take into account risk.)}

\section{Optimal stopping}



\section{Pricing a European option with a binomial model}

\section{Pricing an American option with a binomial model}\label{sec:american}

A European put option allows the holder to decide, at maturity time $T$, whether to exercise the option and collect the payoff $K-S_T$, or to let the contract expire worthless if the strike $K$ is less than the spot price $S_T$. The point is: The European contract stipulates that the holder can \textit{only} exercise at the maturity time $T$, not before. An \textit{American put option}, on the other hand, is identical to its European cousin in every respect, \textit{except} that it allows the holder to exercise at \textit{any} time $t$, up to and including the maturity time $T$. So the decisions facing the owner of an American put option are, at all times $t$ with $0 \leq t < T$, whether to exercise the option and collect
	\[
	\Phi(S_t) \stackrel{\text{def}}{=} \max(K-S_t,0),
	\]
or to keep the contract alive and continue to time $t+1$. The decisions are made in the hopes to maximize profit, of course.

For the sake of argument, let's suppose that the owner of the option has let the contract carry to time $t=T-1$ without exercise, one time period before maturation at time $t=T$. They must then decide whether to exercise immediately and collect $\Phi(S_{T-1})$, or to wait and collect $\Phi(S_T)$. In making this decision, the difficulty is two-fold:
\begin{enumerate}
\item \textit{While $\Phi(S_{T-1})$ is a known quantity at time $t=T-1$, the quantity $\Phi(S_T)$ is not.} The option owner cannot see into the future, even just one time period ahead.
\item Even if $\Phi(S_T)$ were a deterministic quantity, the payoffs occur at different times, one at the present time $t=T-1$ and the other at a future time $t=T$, and so \textit{proper discounting must be applied to obtain the present value of $\Phi(S_T)$.}
\end{enumerate}
We encourage the reader to return to the discussion at the beginning of Section~\ref{sec:claim-pricing}, and notice how these two problems align exactly with the problems we discussed there. The lessons learned in that section suggest that---provided the price process of the stock is a binomial model and the market is arbitrage-free---we can solve the first of these problems by replacing the random quantity $\Phi(S_T)$ with its expectation $E_Q(\Phi(S_T))$ relative to the risk-neutral measure $Q$, and solve the second problem by using the discount factor $1/R$, where $R$ is the risk-free return. Thus, the owner should exercise at time $t=T-1$ if
	\begin{equation}\label{eq:american-intrinsic-continuation}
	\Phi(S_{T-1})(\omega) > \frac{1}{R} E_Q\left(\Phi(S_T)\mid \mathcal{F}_{T-1}\right)(\omega),
	\end{equation}
or defer exercise to time $t=T$ if this inequality is reversed. The quantity $\Phi(S_{T-1})$ on the left-hand side of this inequality is called the \textit{intrinsic value} of the option at time $t=T-1$, while the quantity on the right-hand side is called its \textit{continuation value}.

Suppose, for each $\omega\in \Omega$, that we define
	\[
	V_T = \max{ \left[ K-S_T, 0 \right]}
	\]
and then set
	\[
	V_t = \max{ \left[ K-S_t, \frac{1}{R} E_Q\left( V_{t+1} \mid \mathcal{F}_t \right)\right]}
	\]
for $0 \leq t < T$.


\end{document}